% finding
In this study, we leveraged vertex-wise multimodal MRI from birth to five years of age to chart how structural and functional hierarchies co-develop to shape structure–function coupling in the human cortex. Across this critical window, we observed a pronounced global strengthening of structure–function coupling (SFC) accompanied by parallel increases in intracortical myelination and a shift toward faster, more desynchronized temporal dynamics indexed by the Hurst exponent. Importantly, these changes were not spatially uniform: SFC, structural connectivity (SC), functional connectivity (FC), myelination, and Hurst all exhibited systematic variation along the sensorimotor–association (SA) axis, with unimodal systems showing early, high coupling and association systems displaying delayed and protracted refinement. Together, these findings reveal that the SA hierarchy provides a common developmental scaffold along which cortical wiring, microstructure, and dynamics are jointly tuned, establishing the early architecture that will later support specialized cognition and potentially confer vulnerability to atypical trajectories.

We demonstrate that functional gradient differentiation and excitation-inhibition dynamics are spatially coupled during early infancy, with primary sensory regions showing both the strongest gradient maturation and the largest shifts toward excitation dominance. This coupling suggests that local physiological changes in E/I balance may drive the emergence of macroscale functional organization. The developmental trajectory follows a sensorimotor-first pattern: primary regions differentiate rapidly in the first 6 weeks before stabilizing, while association cortices lag behind, consistent with hierarchical models of cortical maturation. Our vertex-level framework, combining functional gradient analysis with Hurst exponent mapping, enables fine-scale characterization of neurophysiology in the infant brain. Applied to the HBCD cohort, this approach establishes normative trajectories of gradient-E/I coupling that may facilitate early detection of atypical patterns implicated in neurodevelopmental disorders. As longitudinal HBCD data become available, this framework will enable tracking of individual developmental trajectories across infancy.

We further demonstrate that OS-SVD enables high-quality individual gradients that align with population-level patterns. Across denoising benchmarks, OS-SVD provided the strongest improvement in temporal SNR while preserving structural information in residual maps, and it remained robust under temporally undersampled acquisitions. These results indicate that individual gradients can be reliably obtained with under 4 minutes of rs-fMRI acquisition time, supporting faster and more practical infant imaging protocols.

% SC-> SFC, FC -> sFC 
A central implication of our findings is that the unimodal-transmodal gradient provides a common scaffold for the co-development of structural and functional hierarchies that jointly shape SFC. This interpretation is supported by studies using definitions of SFC other than the correlational approach here. For instance, (i) harmonic decomposition of the connectome expresses BOLD signals in graph-frequency space, with structure-function decoupling defined by the ratio of high (decoupled) to low (coupled) frequency components\cite{vazquez-rodriguezGradientsStructureFunction2019,liuTimeresolvedStructurefunctionCoupling2022}, and (ii) in predictive regression frameworks, SFC is indexed by how well a region's empirical FC can be explained from structural features such as Euclidean distance, shortest path length, and communicability derived from the connectome. These complementary methods converge on a heterogeneous pattern of structure-function decoupling along the sensory-association hierarchy, indicating that the phenomenon is robust across analytic definitions of SFC. Notably, it is also robust across developmental stages: prior work in adults and older youth has shown the gradual decoupling of function from structure along the unimodal-transmodal gradient\cite{marguliesSituatingDefaultmodeNetwork2016,fotiadisMyelinationExcitationinhibitionBalance2023,suarezLinkingStructureFunction2020a,guHeritabilityInterindividualVariability2021}, and our vertex-wise analyses extend this principle into the first five years of life, demonstrating that SFC, SC, and FC already obey a shared S-A hierarchy during a period of rapid cortical maturation.

The unimodal-transmodal axis has been proposed as a principal coordinate system that captures the developmental chronology of multiple neurobiological properties, including structural connectivity strength, functional connectivity organization, intracortical myelination, and intrinsic activity timescales\cite{sydnorNeurodevelopmentAssociationCortices2021,xuStructuralConnectivityMatures2024,fotiadisMyelinationExcitationinhibitionBalance2023}. Echoing longitudinal tractography work in older children and adolescents, we observed that unimodal visual and somatomotor systems are both strongly wired and tightly structure-function coupled from early in life, whereas attentional, control, default mode, and limbic networks occupy higher positions along the SA axis, with weaker initial SFC but more protracted strengthening\cite{lebelLongitudinalDevelopmentHuman2011}.  These results suggest that the SA hierarchy not only organizes the mature connectome, but also governs the temporal ordering with which structural constraints and functional specialization are imposed during early stages of development.

% thus, proof of myeliation correlation , myeliation ->SC -> SFC
A second implication of our findings is that hierarchical differences in SFC appear to arise from the coordinated maturation of structural and functional components of cortical circuitry. In our framework, intracortical myelination serves as a proxy for the structural substrate that supports communication, whereas the Hurst exponent indexes local functional dynamics related to excitation-inhibition (E/I) balance.

%H relate with I in one paper, they find H + from 4-5 but they do 400 ROI 
The functional component indicates that heterogeneity in excitation/inhibition balance is inversely related to structure--function coupling. Regions with higher Hurst exponents---reflecting more temporally autocorrelated dynamics and lower E/I ratios---tended to exhibit stronger and less temporally variable SFC. Because the E/I ratio is defined as excitation relative to inhibition, lower E/I values correspond to relatively stronger inhibitory influence; under this interpretation, our observed positive Hurst--SFC relationship supports a direct link between inhibitory-dominant regimes and tighter SC--FC alignment. Developmentally, Hurst values decreased across 0--5 years, indicating a shift toward higher E/I ratios (greater relative excitation), with the most pronounced changes occurring in the first six months of life (Fig.~\ref{fig:hurst_sa}). Notably, both the absolute range and the direction of Hurst changes that we observe are consistent with recent developmental work reporting rapid decreases in scale-free exponents during infancy and early childhood\cite{stanyardAperiodicHurstEEG2024a,houtmanSTXBP1SyndromeCharacterized2021,nishioHurstExponentMarker2024}. We treat Hurst as an index of E/I-related temporal dynamics, but acknowledge that in early development it likely reflects a composite of changes in excitation, inhibition, conduction speed, and global brain state. Taken together with the myelination findings, our results support a view in which SFC in early childhood is jointly shaped by a structural component (intracortical myelination) and a functional component (E/I-related dynamics), whose coordinated but non-uniform maturation underlies the emergent hierarchy of structure--function alignment.
% Consistent with this view, microcircuit studies and modeling work have shown that both excitatory receptor composition and inhibitory interneuron architecture systematically vary along the cortical hierarchy, and that regional differences in Hurst and related 1/f measures covary with inhibitory markers, including parvalbumin (PV$^{+}$) interneuron gene expression\cite{nishioHurstExponentMarker2024}.
% From unimodal sensory and motor areas to heteromodal association cortex, there is a progressive increase in synaptic excitation, supported by greater dendritic complexity and spine density, as well as higher expression of GluN2B-containing NMDA receptors that mediate slow-gating excitatory transmission and support prolonged integration and persistent firing, particularly in prefrontal neurons\cite{wangMacroscopicGradientsSynaptic2020a,FotiadisBiological2026,hensonDevelopmentalRegulationNMDA2008}.
% In contrast, primary sensory regions rely more heavily on fast AMPA receptor-mediated transmission, which is better suited to rapid encoding of sensory inputs and efficient bottom-up signal propagation\cite{wangMacroscopicGradientsSynaptic2020a,YangSynaptic2018}.
% A parallel shift is observed in inhibitory circuitry: in unimodal cortex, inhibition is dominated by PV$^{+}$ parvalbumin-expressing interneurons that target perisomatic compartments and tightly regulate spike output, whereas in association cortex a larger proportion of inhibition is mediated by CB$^{+}$ calbindin- and CR$^{+}$ calretinin-expressing interneurons that innervate distal dendrites and gate the integration of convergent long-range and neuromodulatory inputs\cite{gonzalez-burgosFunctionalMaturationGABA2015,fishParvalbuminContainingChandelierBasket2013,fungExpressionInterneuronMarkers2010}.
% Together, these gradients in excitatory receptor composition and inhibitory targeting provide a mechanistic substrate for the more persistent, integrative dynamics indexed by higher Hurst exponents in association cortex, and help explain why developmental changes in E/I-related dynamics and SFC are coordinated yet heterogeneous along the S--A axis.

At the same time, the Hurst exponent remains an indirect and coarse marker of underlying neurophysiology. HE quantifies the degree of long-range temporal dependence and scale-free structure in the BOLD signal\cite{gaoInferringSynapticExcitation2017}, but it does not in itself decompose activity into specific frequency bands or distinguish the contributions of excitatory versus inhibitory neuronal populations\cite{carterlenoInfantExcitationInhibition2022,stanyardAperiodicHurstEEG2024a}. As such, our interpretation of Hurst in terms of E/I balance rests on converging evidence from biophysical models and gene-expression data rather than a one-to-one mapping between H and particular cell types or spiking regimes\cite{trakoshisIntrinsicExcitationinhibitionImbalance2020,nishioHurstExponentMarker2024}. We therefore view Hurst as a useful macroscopic summary of E/I-related temporal dynamics, while recognizing that it aggregates multiple developmental processes and should be interpreted in concert with microscale and molecular measurements.

Myelination provides a biologically grounded and quantifiable marker of structural maturation, integrating cellular and molecular processes including glial-neuronal interactions as well as genetic and environmental influences\cite{naveMyelinationSupportAxonal2010,hughesGlialCellsPromote2021}. Quantitative MRI studies in infancy and early childhood indicate that intracortical and white-matter myelination follow a sensory-to-association gradient. Occipital, parietal, and sensorimotor pathways show rapid, early increases in myelin content and mature earlier and faster, whereas frontal and temporal cortices exhibit a more protracted, slower myelination trajectory, with frontal white matter in particular showing a longer lag and more prolonged development consistent with late-maturing association systems\cite{deoniCorticalMaturationMyelination2015a,yuInitialMyelinationCentral2023}.
Along the S-A axis, dense intracortical myelin in unimodal sensory and motor cortices has been shown to limit the formation of new axonal projections and synapses, thereby stabilizing existing circuits and reinforcing their functional specialization\cite{deoniCorticalMaturationMyelination2015a,millerProlongedMyelinationHuman2012,FaccaMultiscale2024}. In this context, high levels of intracortical myelin in unimodal regions may facilitate the transmission of functional signals that remain closely aligned with underlying anatomical pathways. Consistent with this interpretation, the unimodal pole showed larger age-related increases and higher overall myelin levels, which were in turn associated with stronger SFC and steeper SFC maturation (Fig.~\ref{fig:myelin_sa}). In contrast, transmodal cortex exhibits comparatively sparse intracortical myelination, preserving greater capacity for synaptic reorganization and flexible, context-dependent dynamics that can diverge from the structural scaffold. Such reduced myelin constraint may help explain weaker structure–function coupling in these regions, while simultaneously positioning them to support adaptive behavior and learning.

%reason why increase in inhibition right decide this is not useful i cut it off
% Transcriptomic and proteomic studies of the human cortex have shown that there's change in the expression of excitatory and inhibitory genes, such as NMDA receptor subunits, calcium-binding proteins that are expressed in different types of interneurons, and GABAa receptor subunits\cite{gonzalez-burgosFunctionalMaturationGABA2015,fishParvalbuminContainingChandelierBasket2013,fungExpressionInterneuronMarkers2010}. Molecular shifts mirror changes in neuronal functional properties. 
% In prefrontal cortex, there's increases in the subunit composition fo GABAa receptors, from a2 to a1 containing receptors, which have a faster decay time, resulting in faster synaptic inhibition\cite{hensonDevelopmentalRegulationNMDA2008}. 

% say something about e/i balance and e/i ratio is different 
